% ===== PRÉAMBULE CANONIQUE — à coller VERBATIM en tête de CHAQUE .tex =====
\documentclass[11pt,a4paper]{article}
\usepackage[utf8]{inputenc}\usepackage[T1]{fontenc}\usepackage{lmodern}
\usepackage[french]{babel}
\usepackage{amsmath,amssymb,mathtools}\usepackage{icomma}
\usepackage[version=4]{mhchem}          % \ce{...} : équations & formules
\usepackage{siunitx}\sisetup{locale=FR} % \qty{}{}, \si{}, \ang{}
\usepackage{chemfig}                    % \chemfig{...} : structures développées
\usepackage{xcolor}\usepackage{colortbl}
\usepackage{tikz}\usepackage{pgfplots}\pgfplotsset{compat=1.18}
\usepackage[most]{tcolorbox}\usepackage{enumitem}\usepackage{array,booktabs}
\usepackage[a4paper,margin=2.2cm]{geometry}
\usetikzlibrary{arrows.meta,calc,angles,quotes,positioning,patterns,backgrounds,shapes.geometric}
\definecolor{primary}{RGB}{222,49,99}\definecolor{primarydark}{RGB}{150,22,55}
\definecolor{defcol}{RGB}{0,114,178}\definecolor{methcol}{RGB}{52,92,148}
\definecolor{excol}{RGB}{0,135,90}\definecolor{attncol}{RGB}{200,40,55}
\tcbset{boxbase/.style={enhanced,breakable,boxrule=0.9pt,arc=2.5pt,
  left=3mm,right=3mm,top=2.3mm,bottom=2.3mm,fonttitle=\bfseries,coltitle=white}}
% --- boîtes TUTORIEL (noms EXACTS, ne pas renommer) ---
\newtcolorbox{cle}[1][]{boxbase,colback=defcol!6,colframe=defcol,
  title={\textsc{À retenir}\ifx\relax#1\relax\relax\else\textmd{ : \itshape #1}\fi}}
\newtcolorbox{methode}[1][]{boxbase,colback=methcol!7,colframe=methcol,sharp corners,
  title={\textsc{Méthode}\ifx\relax#1\relax\relax\else\textmd{ --- \itshape #1}\fi}}
\newtcolorbox{exemplebox}[1][]{boxbase,colback=excol!5,colframe=excol,
  title={\textsc{Exemple}\ifx\relax#1\relax\relax\else\textmd{ : \itshape #1}\fi}}
\newtcolorbox{piege}[1][]{boxbase,colback=attncol!6,colframe=attncol,
  title={\textsc{Piège}\ifx\relax#1\relax\relax\else\textmd{ : \itshape #1}\fi}}
\newtcolorbox{remarque}[1][]{boxbase,colback=black!4,colframe=black!45,
  title={\textsc{Remarque}\ifx\relax#1\relax\relax\else\textmd{ : \itshape #1}\fi}}
% --- boîtes EXERCICE / CORRIGÉ (noms EXACTS, auto-comptées) ---
\newtcolorbox[auto counter]{exo}[1][]{boxbase,colback=primary!4,colframe=primary,
  title={\textsc{Exercice}~\thetcbcounter\ifx\relax#1\relax\relax\else\textmd{ --- \itshape #1}\fi}}
\newtcolorbox[auto counter]{corrige}[1][]{boxbase,colback=excol!4,colframe=excol,
  title={\textsc{Corrigé}~\thetcbcounter\ifx\relax#1\relax\relax\else\textmd{ --- \itshape #1}\fi}}
\newcommand{\R}{\mathbb{R}}\newcommand{\N}{\mathbb{N}}\newcommand{\Z}{\mathbb{Z}}\newcommand{\ds}{\displaystyle}
% (un \usepackage{fancyhdr}+en-tête est possible ; il est ignoré côté web)
% =========================================================================
\begin{document}
\sisetup{per-mode=symbol}
\begin{center}\small\itshape On rappelle : $\log 2 \approx 0{,}30$,\quad $\log 3 \approx 0{,}48$.\end{center}

\begin{exo}[choisir et doser un couple tampon]
On veut préparer un tampon de $\mathrm{pH} = 5{,}0$. On dispose des couples : \ce{HCOOH/HCOO-} ($3{,}8$),
\ce{CH3COOH/CH3COO-} ($4{,}8$), \ce{H2CO3/HCO3-} ($6{,}4$), \ce{NH4+/NH3} ($9{,}2$).
\begin{enumerate}[label=\alph*)]
  \item Quel couple choisir ? Justifie (zone d'efficacité $\mathrm{p}K_a \pm 1$).
  \item Calcule le rapport $\dfrac{[\ce{A-}]}{[\ce{HA}]}$ nécessaire.
\end{enumerate}
\end{exo}

\begin{exo}[le tampon à la demi-équivalence]
On part de \qty{100}{mL} d'acide acétique \ce{CH3COOH} à \qty{0,20}{mol\per L} ($\mathrm{p}K_a = 4{,}8$), que
l'on neutralise progressivement par \ce{NaOH} à \qty{0,20}{mol\per L}.
\begin{enumerate}[label=\alph*)]
  \item Quel volume de \ce{NaOH} faut-il verser pour que le mélange soit un tampon de $\mathrm{pH} = 4{,}8$ ?
  \item Quel nom porte ce point du titrage ?
\end{enumerate}
\end{exo}

\begin{exo}[titrage par un diacide]
On neutralise \qty{30}{mL} de \ce{NaOH} à \qty{0,10}{mol\per L} par de l'acide sulfurique \ce{H2SO4}
(diacide) à \qty{0,050}{mol\per L}.
\begin{enumerate}[label=\alph*)]
  \item Écris la relation à l'équivalence en tenant compte des \textbf{deux} acidités de \ce{H2SO4}.
  \item Calcule le volume de \ce{H2SO4} à l'équivalence.
\end{enumerate}
\end{exo}

\begin{exo}[un tampon encaisse un ajout d'acide fort]
On dispose de \qty{100}{mL} d'un tampon contenant \ce{HCOOH} à \qty{0,50}{mol\per L} et \ce{HCOONa} à
\qty{0,50}{mol\per L} ($\mathrm{p}K_a = 3{,}8$). On ajoute \qty{10}{mL} de \ce{HCl} à \qty{1,0}{mol\per L}.
\begin{enumerate}[label=\alph*)]
  \item Écris la réaction de l'acide fort avec le tampon, et calcule les nouvelles quantités de \ce{HCOOH} et
        \ce{HCOO-}.
  \item En déduis le pH final (Henderson).
\end{enumerate}
\end{exo}

\begin{exo}[quel sel donne une solution basique ?]
On donne : \ce{HCN} ($\mathrm{p}K_a = 9{,}2$) est un acide faible ; \ce{HBr}, \ce{HNO3}, \ce{HCl} sont des
acides forts ; \ce{KOH}, \ce{Ca(OH)2} sont des bases fortes.
\begin{enumerate}[label=\alph*)]
  \item Pour chacun des sels \ce{Ca(NO3)2}, \ce{KCN}, \ce{NaCl}, \ce{NH4Br}, indique de quel acide et de quelle
        base il provient.
  \item Lequel donne une solution \textbf{basique} ?
\end{enumerate}
\end{exo}

\end{document}
